% Generated draft: requires mathematical and source review.
\documentclass[11pt]{article}
\usepackage[margin=25mm]{geometry}
\usepackage{fontspec}
\setmainfont{Latin Modern Roman}
\usepackage{amsmath,amssymb,mathtools,unicode-math}
\setmathfont{Latin Modern Math}
\usepackage{enumitem,needspace,xurl,hyperref}
\setlength{\parindent}{0pt}
\setlength{\parskip}{5pt}
\setlength{\emergencystretch}{3em}
\newcommand{\N}{\mathbb{N}}
\newcommand{\Z}{\mathbb{Z}}
\newcommand{\Q}{\mathbb{Q}}
\newcommand{\R}{\mathbb{R}}
\newcommand{\C}{\mathbb{C}}
\newcommand{\F}{\mathbb{F}}
\newcommand{\A}{\mathbb{A}}
\newcommand{\PP}{\mathbb{P}}
\newcommand{\E}{\mathbb{E}}
\newcommand{\eps}{\varepsilon}
\newcommand{\dd}{\,\mathrm{d}}
\begin{document}
\section*{UnsolvedMath ID: 800787}
\textbf{Source problem code:} catalogue entry
\subsection{Definitions and mathematical statement}
Use the standard germ conventions as follows. Let \(n\geq 1\) be an integer and let \(O=0\in\C^n\). Write \(\operatorname{End}(\C^n,O)\) for the holomorphic map germs \(f:(\C^n,O)\to(\C^n,O)\) fixing \(O\). A germ is non-invertible if it has no holomorphic inverse germ, equivalently \(\det Df(O)=0\); it is dominant if \(\det Df\) is not identically zero as a germ.

For a nonzero scalar holomorphic germ \(h\), let \(\operatorname{ord}_O(h)\) be the least total degree of a nonzero term of its Taylor series, and put \(\operatorname{ord}_O(0)=+\infty\). For \(f=(f_1,\ldots,f_n)\), define its order and asymptotic order by
\[
 c(f)=\min_{1\leq i\leq n}\operatorname{ord}_O(f_i),
 \qquad
 c_\infty(f)=\lim_{k\to\infty}c(f^{\circ k})^{1/k},
\]
where \(f^{\circ k}\) is the \(k\)-fold compositional iterate. These are the standard interpretations of the notation, which the supplied excerpt does not itself define.

\textbf{Problem OP5 (Abate, 2011).} For every non-invertible dominant germ \(f\in\operatorname{End}(\C^n,O)\), ask whether both of the following hold:
\begin{enumerate}
\item The number \(c_\infty=c_\infty(f)\) is an algebraic integer of algebraic degree at most \(n\); equivalently, there exist integers \(a_0,\ldots,a_{n-1}\) such that
\[
 c_\infty^n+a_{n-1}c_\infty^{n-1}+\cdots+a_1c_\infty+a_0=0.
\]
\item There exists a real number \(\delta=\delta(f)\in(0,1]\) such that, for every integer \(k\geq1\),
\[
 \delta\,c_\infty(f)^k\leq c(f^{\circ k})\leq c_\infty(f)^k.
\]
\end{enumerate}
The source calls the algebraic degree the ``order'' of the algebraic integer. Its displayed polynomial begins with \(c_\infty^n\) but then uses the index \(a_{k-1}c_\infty^{k-1}\), with no corresponding \(k\) specified there. The display above normalizes that index to \(n-1\), following the stated degree bound and coefficient list; this editorial interpretation requires review.

\textbf{Supplied context.} The source compares the order of a germ, as a measure of attraction toward the origin, with polynomial degree, as a measure of attraction toward infinity. It attributes to Favre and Jonsson (2011), reference [FJ2], the assertion that for every polynomial map \(F:\C^2\to\C^2\), the sequence of degrees of its iterates satisfies a linear recursion with integer coefficients. This is supplied background, not an independently verified theorem here. The excerpt then asks generally about an analogous recursion for orders of iterates of germs, but supplies no complete additional formulation; that trailing question is not incorporated into OP5.
\subsection{Short English statement}
Abate's OP5 asks whether every non-invertible dominant holomorphic germ fixing the origin in \(\C^n\) has asymptotic order that is an algebraic integer of degree at most \(n\), and whether the orders of all its positive iterates stay between \(\delta c_\infty^k\) and \(c_\infty^k\) for one constant \(\delta\in(0,1]\) depending on the germ. Standard definitions and an explicitly identified correction to a coefficient index are used. The registry labels the problem open; that status is supplied metadata.
\subsection{Sources}
\begin{itemize}
\item Local UnsolvedMath registry: permanent ID \(800787\), source problem code \texttt{catalogue entry}, title ``Local Discrete Holomorphic Dynamics --- Open Problem OP5.'' The supplied record attributes the problem to Abate (2011) and labels its status \texttt{open}; this label is not independently verified. Upstream revision: \texttt{d371ee917ac9647fbb06cd6fc529deec6f3dda82}.
\item Abate, \emph{Open Problems in Local Discrete Holomorphic Dynamics} (2011), Problem OP5, \url{https://arxiv.org/abs/1107.4876}. The supplied provenance records access on 2026-10-16 and extraction by \texttt{arxiv\_latex}; these are supplied provenance claims, not a claim of access or verification for this draft.
\item Favre and Jonsson (2011), reference [FJ2] in the supplied excerpt: attributed background concerning integer linear recursions for degrees of iterates of polynomial maps \(\C^2\to\C^2\). No full bibliographic entry or separate URL for [FJ2] is supplied.
\end{itemize}
\end{document}
